% Generated draft: requires mathematical and source review.
\documentclass[11pt]{article}
\usepackage[margin=25mm]{geometry}
\usepackage{fontspec}
\setmainfont{Latin Modern Roman}
\usepackage{amsmath,amssymb,mathtools,unicode-math}
\setmathfont{Latin Modern Math}
\usepackage{enumitem,needspace,xurl,hyperref}
\setlength{\parindent}{0pt}
\setlength{\parskip}{5pt}
\setlength{\emergencystretch}{3em}
\newcommand{\N}{\mathbb{N}}
\newcommand{\Z}{\mathbb{Z}}
\newcommand{\Q}{\mathbb{Q}}
\newcommand{\R}{\mathbb{R}}
\newcommand{\C}{\mathbb{C}}
\newcommand{\F}{\mathbb{F}}
\newcommand{\A}{\mathbb{A}}
\newcommand{\PP}{\mathbb{P}}
\newcommand{\E}{\mathbb{E}}
\newcommand{\eps}{\varepsilon}
\newcommand{\dd}{\,\mathrm{d}}
\begin{document}
\section*{UnsolvedMath ID: 30002378}
\textbf{Source problem code:} catalogue entry
\subsection{Definitions and mathematical statement}
Let \(M\) denote a compact surface with the real-analytic structure implicit in the source, and let \(\mu\) denote its prescribed area datum. The source does not specify whether boundary is allowed, whether connectedness or orientability is required, or the regularity and precise meaning of \(\mu\). These conventions remain unresolved. When \(\mu\) is an area form, the intended group is
\[
\operatorname{Symp}^{\omega}_{\mu}(M)
=\{f:M\to M\mid f\text{ is a real-analytic diffeomorphism and }f^{*}\mu=\mu\},
\]
with composition as its group operation. If \(\mu\) instead denotes an area measure, the preservation condition is \(\mu(f^{-1}(A))=\mu(A)\) for every measurable set \(A\). These interpretations must be distinguished, particularly concerning orientability and orientation-reversing maps; the supplied record does not settle the convention.

A group \(G\) is finitely generated if \(G=\langle S\rangle\) for some finite subset \(S\subseteq G\). For a group \(H\), set \(H^{(0)}=H\) and
\[
H^{(n+1)}=\left\langle xyx^{-1}y^{-1}:x,y\in H^{(n)}\right\rangle.
\]
The group \(H\) is solvable if \(H^{(r)}=\{e\}\) for some integer \(r\geq 0\). A group \(G\) is virtually solvable if it has a solvable subgroup \(H\leq G\) of finite index. Let \(F_2=\langle a,b\mid\ \rangle\) be the free group on two generators; containing a nonabelian free group is equivalent to admitting an injective homomorphism \(F_2\hookrightarrow G\).

\textbf{Conjecture (supplied surface formulation).} For every compact surface \(M\), with the analytic structure and preserved area datum \(\mu\) understood as above, and every finitely generated subgroup \(G\leq\operatorname{Symp}^{\omega}_{\mu}(M)\),
\[
\bigl(\exists H\leq G:\ [G:H]<\infty\ \text{and}\ \exists r\geq 0,\ H^{(r)}=\{e\}\bigr)
\quad\text{or}\quad
\bigl(\exists\text{ an injective homomorphism }F_2\to G\bigr).
\]
The source leaves the admissible choices and quantification of \(\mu\), as well as the surface conventions, implicit. Thus this is a statement-only transcription requiring those qualifications to be resolved before a fully specified formulation is accepted. The accompanying literature assessment concerns circle diffeomorphisms, whereas this conjecture concerns area-preserving surface diffeomorphisms; that assessment does not establish the literature status of this surface formulation.
\subsection{Short English statement}
The supplied conjecture asserts that every finitely generated group of real-analytic area-preserving diffeomorphisms of a compact surface has a solvable subgroup of finite index or contains a free group on two generators. Surface and area-preservation conventions remain unspecified. The supplied open label and verification claims are curation metadata; the accompanying circle-diffeomorphism literature assessment does not verify this surface conjecture.
\subsection{Sources}
\begin{itemize}
\item Local UnsolvedMath registry: supplied complete source record, permanent ID \(30002378\), exact source problem code \texttt{catalogue entry}, in the Oberwolfach Reports Open Problems collection. The record supplies proposed year 2013 and no individual proposer.
\item Supplied source citation: \emph{Dynamische Systeme (2014)}, Oberwolfach Reports, DOI \texttt{10.4171/owr/2013/34}, pp.\ 1975--2033. \url{https://doi.org/10.4171/owr/2013/34}. The background also identifies this as Oberwolfach Report 34/2013; the differing year labels are retained without reconciliation.
\item Supplied literature source: \emph{Probabilistic Tits alternative for circle diffeomorphisms}. \url{https://arxiv.org/abs/2412.08779}.
\item Literature checked at: \texttt{2026-08-21}. This is supplied dated curation, not independently verified. The supplied assessment says that the full Tits-alternative assertion for finitely generated subgroups of real-analytic circle diffeomorphisms remains unresolved and that later results give weaker dynamical or probabilistic alternatives. It explicitly calls for independent verification of time-sensitive status and cited claims. Its applicability to the surface conjecture is unresolved.
\item The registry labels the problem \texttt{open} and its statement \texttt{corrected\_verified}, and reports that the mathematical content was checked against the official Oberwolfach report with notation, context, and question form normalized. These are supplied verification claims, not an independent source check performed for this draft.
\end{itemize}
\end{document}
